\documentclass[a4j,dvipdfmx]{jsarticle}
\usepackage{amsmath,amssymb}
\usepackage{siunitx}
\usepackage{ascmac}
\usepackage{amsthm}
\usepackage{bm}

% \usepackage[margin=10truemm]{geometry}

\newcommand{\R}{\mathbb{R}}
\newcommand{\rank}{\mathrm{rank}}
\newcommand{\tr}{\mathrm{tr}}
\newcommand{\Ker}{\mathrm{Ker}\hspace{0.5mm}}
\newcommand{\E}{\mathcal{E}}
\newcommand{\F}{\mathcal{F}}
\renewcommand{\labelenumi}{(\arabic{enumi})}
\renewcommand{\Im}{\mathrm{Im}\hspace{0.5mm}}
% \usepackage[dvipdfmx]{hyperref}
\begin{document}
    \part*{第四回 線型代数 問題}
    \subsection*{問題1}
        $n$次正方行列$A$に対し, 適当な(複素)正則行列$P$を取って, $P^{-1}AP$が対角行列になるようにできる(=対角化可能である)必要十分条件を答えよ.
    \subsection*{問題2}
        線形変換$T$の相異なる固有値に対する固有ベクトルは互いに線形独立であることを示せ.
    \subsection*{問題3}
        次の行列が対角化可能なら対角化し, そうでなければ理由を述べよ.
        \begin{equation*}
            (1)\begin{pmatrix}
                6 & -3 & -7\\
                -1 & 2 & 1\\
                5 & -3 & -6
            \end{pmatrix}\quad 
            (2)\begin{pmatrix}
                1  & 2  & 1 \\
                -1 & 4  & 1 \\
                2  & -4 & 0 
            \end{pmatrix}\quad
            (3)\begin{pmatrix}
                1  & 3  & -2 \\
                -3 & 13  & -7 \\
                -5  & 19 & -10 
            \end{pmatrix}
        \end{equation*}
    \subsection*{問題4}
        $A=\begin{pmatrix}1 & 2 & 1\\-1 & 4 & 1\\2 & -4 & 0\end{pmatrix}$とするとき, $A^n\hspace{1mm}(n\in \mathbb{N})$を求めよ.
    \subsection*{問題5}
        漸化式$x_{n+2}=x_{n+1}+x_{n}$を解け.
    \subsection*{問題6}
        微分方程式$y'''-4y''+y'+6y=0$をとけ.
    \subsection*{問題7}
        $A=\begin{pmatrix}1 & 2\\2 & 4\end{pmatrix}$による線形変換を$f:\R^2\to\R^2$とする. 原点を通る直線$L$を$f$で写すと再び$L$に重なる, つまり$L=f(L)$のとき, 直線$L$と直線$y=3$との交点を求めよ.

\end{document}